\paragraph{De l'origine incidencielle à la modification réticulaire.}
\label{inc-ampl-reticulo}
Le drapeau marqué détermine une position au moyen d'une opération
ensembliste équivariante :
\[
 o_*=(H_e\cap H_\varphi)
       \setminus(H^-_\alpha\cup H_{\rm APP})=\{1\}.
\]
Sa fonction consiste à distinguer une composante parmi les douze qui
interviennent dans la réalisation réticulaire. Le code ternaire complet
agit comme code de recollement dans le module discriminant
$(A_2^*/A_2)^{12}$. Sa préimage définit $N(C_W)$, et
l'autodualité, les poids et la distance du code se traduisent,
respectivement, en propriétés d'intégralité et de déterminant, de parité
et en bornes de norme des classes de recollement. L'information symbolique
du code remplit ici une fonction que le support isolé ne remplit plus.

La métrique $A_2$ et la chambre radiale positive sont fixées dans le développement
réticulaire. À l'intérieur de cette chambre, la congruence nécessaire pour un voisin
pair d'indice trois possède douze réalisations de norme minimale $54$,
permutées par le changement de composante distinguée. L'origine $o_*$
sélectionne
\[
 v_\alpha=4\rho_1+\rho_2+\cdots+\rho_{12},\qquad
 v_\alpha^2=2(4^2+11)=54.
\]
Le coefficient $4$ appartient à cette solution de la congruence dans la
chambre déclarée. L'indice $\alpha$ enregistre la provenance
incidencielle de la sélection : son contenu est le repère qui est intervenu dans
la clôture et qui individualise à présent une direction réticulaire.

L'appariement avec $v_\alpha$ détermine le sous-réseau
$N_{\alpha,0}$ d'indice trois, et l'adjonction de la classe
$v_\alpha/3$ produit le voisin $N_\alpha$ de
\eqref{exc:vecino}. Cette modification conserve le rang et rétablit
l'unimodularité et la parité ; la sélection par congruence élimine les racines
antérieures, et les bornes locales excluent les racines dans les nouvelles classes.
L'identification à Leech intervient après ces vérifications. La
réalisation binaire des cinq régions et cette construction réticulaire
sont deux prolongements du repère commun : la seconde reçoit directement
le code ternaire comme donnée de recollement.

